\subsection*{Interview-Taking Advice}

\begin{itemize}
	\item Relax. Your interview will be between several engineering/physics/maths nerds talking about engineering/physics/maths problems - enjoy it.
	\item Not being able to solve a problem but instead exploring multiple incorrect paths through a problem that ended up being wrong would be a better interview performance than simply writing down the correct answer. The important part is to motivate why you decided on those paths and why you gave up on those paths. This is not a test of whether you can get questions correct - there are entrance exams for this purpose.
	\item Interviewers are also looking for how well they can work with a student in a tutorial and want to hear your thought process, even if you are going down the wrong path or saying something incorrect. This is a natural part of exploring a problem and does not count against you.
	\item Can you work with a tutor? The interview is more collaborative than combative and the interviewer will be trying to help you when they see you are stuck - make sure to use any hints.
	\item You can and should ask clarifying questions\footnote{The interviewer may have deliberately omitted an assumption that changes the question such as whether a function is real-valued, whether friction is present et cetera. You may want to discuss the suitability of making a given assumption and then ask the interviewer whether that it ok once you have decided.}, but avoid asking questions that try to probe the route through the problem - this part is meant to come from you.
	\item If you spot a pattern that ends up being useful, the interviewer may want you to explain/prove that later. It may be helpful to your problem solving process to note that your conjecture is unproven but assume it is true for the time being.
	\item This is also a test of problem solving. Try to demonstrate that you can build intuition for an unseen problem, break a problem down into simpler parts, and apply understanding gained from previously solved problems.
	\item A levels do not prepare you well to make formal arguments. The interviewers will not be expecting you to be amazing at this as it is something you will be trained in at university, but it is something to begin thinking about. In an interview they will still expect you to collate your fuzzy intuitions about what a solution may look like into a coherent argument however.
\end{itemize}

\subsection*{Problem-Solving Advice}

\begin{itemize}
	\item If the original problem is too hard, consider a simplified version. For example if the problem is very abstract, looking at a concrete example may help.
	\item Make observations about the problem before diving in. For example when deriving the volume of a sphere of radius $r$, you expect it to be proportional to $r^3$. Any simple bounds or relationships you spot can also help.
	\item Graphs are your friend. For example, proving that a cubic having three roots is equivalent to having two stationary points with opposite sign.
	\item Have you solved a similar problem before which may provide useful intuition?
	\item Look for simplifications such as symmetries. For example, when finding the average distance between two randomly chosen points on the circumference of a circle, the circle can always be oriented such that one of the point is at the top of the circle.
	\item If your answer involves considering a large number of cases then you are probably doing something wrong. You should typically only begin listing out cases to try and spot a pattern, but do not intend on finishing.
	\item Considering how much information you need to extract and how much information you can extract is often useful when finding strategies. For example, 3 “yes/no” questions can only distinguish between 8 states.
	\item Understand the pigeonhole principle: \url{https://www.youtube.com/watch?v=TCZ3YwbcDaw}.
	\item Understand how possibilities combine. If $A$ has $n$ possibilities and $B$ has $m$ possibilities then $A$ and $B$ together have at most $nm$ possibilities. Ask your teacher for a probability exercise comparing the configuration space for independent and non-independent events to understand these ideas further.
\end{itemize}